معادله دیفرانسیل داده‌شده به صورت زیر است:
\begin{equation*}
y'(t) + a \cdot y(t) = b \cdot x(t)
\end{equation*}

\subsectionaddtolist{الف}
با اعمال تبدیل فوریه به دو طرف معادله دیفرانسیل و استفاده از خاصیت مشتق‌گیری داریم:
\begin{equation*}
j\omega Y(\omega) + a Y(\omega) = b X(\omega) \implies (j\omega + a) Y(\omega) = b X(\omega)
\end{equation*}
تابع تبدیل سیستم $H(\omega)$ به صورت زیر تعریف می‌شود:
\begin{equation*}
H(\omega) = \frac{Y(\omega)}{X(\omega)} = \frac{b}{j\omega + a}
\end{equation*}
با گرفتن تبدیل فوریه معکوس، پاسخ ضربه در حوزه زمان (با فرض $a > 0$) برابر است با:
\begin{equation*}
h(t) = \mathcal{F}^{-1}\left\{ \frac{b}{j\omega + a} \right\} = b \cdot e^{-at} u(t)
\end{equation*}

\subsectionaddtolist{ب}
پاسخ ضربه در حوزه فرکانس همان تابع تبدیل $H(\omega)$ است:
\begin{equation*}
H(\omega) = \frac{b}{j\omega + a}
\end{equation*}
 {اندازه:}
    \begin{equation*}
    |H(\omega)| = \frac{|b|}{\sqrt{\omega^2 + a^2}}
    \end{equation*}
    {فاز:}
    \begin{equation*}
    \angle H(\omega) = \angle b - \tan^{-1}\left(\frac{\omega}{a}\right)
    \end{equation*}
    (با فرض $b>0$ فاز برابر با $-\tan^{-1}(\frac{\omega}{a})$ خواهد بود).

\subsectionaddtolist{ج}
پاسخ پله از انتگرال‌گیری پاسخ ضربه در حوزه زمان به دست می‌آید:
\begin{equation*}
s(t) = \int_{0}^{t} b \cdot e^{-a\tau} d\tau = \left[ -\frac{b}{a} e^{-a\tau} \right]_{0}^{t} = \frac{b}{a} \left(1 - e^{-at}\right) u(t)
\end{equation*}
بررسی حد در بی‌نهایت:
\begin{equation*}
\lim_{t \to \infty} s(t) = \lim_{t \to \infty} \frac{b}{a} \left(1 - e^{-at}\right) = \frac{b}{a}
\end{equation*}

\subsectionaddtolist{د}
طبق قانون ولتاژ کیرشهف در مدار RC متوالی سری:
\begin{equation*}
x(t) = R \cdot i(t) + y(t)
\end{equation*}
از آنجا که جریان خازن برابر است با $i(t) = C \frac{dy(t)}{dt}$، داریم:
\begin{equation*}
x(t) = RC \frac{dy(t)}{dt} + y(t)
\end{equation*}
با تقسیم طرفین بر $RC$ جهت هم‌فرم کردن با معادله اصلی:
\begin{equation*}
y'(t) + \frac{1}{RC} y(t) = \frac{1}{RC} x(t)
\end{equation*}
با مقایسه این رابطه با معادله دیفرانسیل صورت سؤال، مقادیر $a$ و $b$ به دست می‌آیند:
\begin{equation*}
a = \frac{1}{RC} , \quad b = \frac{1}{RC}
\end{equation*}
